\section{Cryptography and Blockchain Foundations}

Blockchain is not a single new technology; it is an architecture built by
combining three foundational cryptographic primitives.

\subsection{Primitive 1: Encryption and Decryption}
Encryption secures communication between parties who may have never met,
ensuring unauthorized eavesdroppers cannot read the data.
\begin{itemize}
    \item \textbf{Symmetric Encryption:} The same key is used to both encrypt and decrypt (e.g.,
          Caesar cipher).
          \begin{itemize}
              \item \textit{Fatal Flaw:} The key distribution problem. Both parties must securely agree on
                    the key beforehand without it being intercepted.
          \end{itemize}
    \item \textbf{Asymmetric Encryption (Public Key Cryptography):} Uses two mathematically
          linked keys.
          \begin{itemize}
              \item \textbf{Public Key:} Broadcasted freely; used by others to \textit{encrypt} messages sent to you.
              \item \textbf{Private Key:} Kept entirely secret; used by you to \textit{decrypt} those messages.
              \item \textit{The Math:} Relies on a ''trapdoor'' function (like ECDSA) that is trivial to
                    compute in one direction but computationally infeasible to reverse (e.g., the
                    Bitcoin keyspace is $2^{256}$).
          \end{itemize}
\end{itemize}

\subsection{Primitive 2: Cryptographic Hashing}
Hashing turns any message (of any length) into a fixed-length fingerprint. It
is used to prove you know something without revealing it (e.g., committing to a
prediction).
\begin{itemize}
    \item \textbf{5 Key Properties of a Cryptographic Hash:}
          \begin{enumerate}
              \item \textbf{Deterministic:} The same input always produces the exact same output.
              \item \textbf{Computable:} Easy to compute forward.
              \item \textbf{Pre-image Resistance:} Hard to reverse (one-way function).
              \item \textbf{Avalanche Effect:} A small change in input completely changes the output.
              \item \textbf{Collision Resistance:} Infeasible to find two different inputs that produce the
                    same output.
          \end{enumerate}
    \item \textbf{Algorithms:} SHA-256 (industry standard, 64 hex characters) vs. MD5 (broken and
          deprecated, as collisions can be generated in seconds).
\end{itemize}

\subsection{Primitive 3: Digital Signatures}
Digital signatures prove that you authored a message and that it hasn't been
altered, completely removing the need for a trusted third party (like a
notary).
\begin{itemize}
    \item \textbf{Sign:} Message + Private Key $\rightarrow$ Digital Signature.
    \item \textbf{Verify:} Message + Digital Signature + Public Key $\rightarrow$ Valid /
          Invalid.
    \item \textbf{The Bitcoin Shift:} In traditional banking, an account exists because a bank's
          database says so. In Bitcoin, ownership is proven mathematically. You control
          an address solely because you can produce a valid digital signature for it.
\end{itemize}

\subsection{The Synthesis: Blockchain Architecture}
When combined, these three primitives allow a network of untrusted parties to
transact with mathematical certainty:
\begin{itemize}
    \item \textbf{Identity:} Asymmetric keys establish identity (Public Key = Address).
    \item \textbf{Authorization:} Digital signatures authorize transactions.
    \item \textbf{Immutability:} Hashing links blocks together. Any tampering alters the hash,
          breaking all subsequent links in the chain.
\end{itemize}

\subsection{Implementation: When to Use Blockchain}
Blockchain is not a universal solution for digitization. A well-designed
centralized database is usually cheaper and faster unless the scenario passes
the \textbf{Boaventura Test}. Consider blockchain if you answer ''yes'' to most of the
following:
\begin{enumerate}
    \item Is the process predominantly cross-organizational?
    \item Are there cross-system discrepancies slowing the business?
    \item Is trust incomplete among transacting parties?
    \item Are intermediaries adding excessive cost, risk, or delay?
    \item Does the current system require periodic offline reconciliation?
    \item Is there a critical need for an immutable audit trail?
\end{enumerate}
\textit{Note on Vulnerabilities:} The cryptography itself (like SHA-256) rarely fails.
Breaches almost always occur at the human layer: poor key management, phishing,
social engineering, or logic bugs written into smart contracts.